% lemniro: {"kind":"note","description":"A circle, its vertical projection, and the differential equation behind a sine wave, with three independent interactive figures.","topic":"Analysis","series":"Geometry in motion","order":1,"date":"2026-10-03","readingMinutes":7,"prerequisites":["Trigonometric functions","The chain rule"],"related":[]}
\documentclass[11pt]{article}
\input{tex/lemniro-preamble}
\title{A circle becomes a sine wave}
\begin{document}
\maketitle

\section{Keep one coordinate}
Fix a period $T>0$ and a phase $\phi\in\R$. A point moving uniformly
around the unit circle has position
\begin{equation}\label{eq:circle}
  p_\phi(t)=\bigl(\cos\theta(t),\sin\theta(t)\bigr),
  \qquad \theta(t)=\frac{2\pi}{T}t+\phi.
\end{equation}
The equation $\cos^2\theta+\sin^2\theta=1$ keeps the point on the circle.
Its second coordinate is the function
\[
  u_\phi(t)=\sin\!\left(\frac{2\pi}{T}t+\phi\right).
\]
Recording that height against time produces a sine wave. The horizontal
coordinate of the graph is time, not the moving point's horizontal coordinate.
The two diagrams therefore share their vertical value while using different
horizontal coordinates.

\LemniroScene{circle-and-sine}

Use the time slider to compare the same height on both diagrams. The moving
dashed connector makes this equality visible. The figure uses an eight-second
cycle and screen coordinates; its vertical pixel coordinate decreases when
the mathematical height increases. That display convention does not change
the function in Equation~\ref{eq:circle}.

\section{The equation behind the motion}
\begin{theorem}[Harmonic motion]\label{thm:harmonic}
Let $\omega=2\pi/T$. For every phase $\phi$, the function
$u_\phi(t)=\sin(\omega t+\phi)$ satisfies
\begin{equation}\label{eq:oscillator}
  u_\phi''(t)+\omega^2u_\phi(t)=0
\end{equation}
and $u_\phi(t+T)=u_\phi(t)$ for every real $t$.
\end{theorem}
\begin{proof}
The chain rule gives
\[
  u_\phi'(t)=\omega\cos(\omega t+\phi),
  \qquad
  u_\phi''(t)=-\omega^2\sin(\omega t+\phi)=-\omega^2u_\phi(t).
\]
Moving the final term to the left proves Equation~\ref{eq:oscillator}.
For periodicity, substitute $t+T$ and use $\omega T=2\pi$:
\[
  u_\phi(t+T)=\sin(\omega t+\phi+2\pi)
             =\sin(\omega t+\phi)=u_\phi(t).
\]
Both claims hold for every real $t$, as required.
\end{proof}

\LemniroScene{circle-and-sine}

Pause near a maximum. The height is momentarily stationary, so its first
derivative is zero, while its second derivative points back toward height
zero. Theorem~\ref{thm:harmonic} states that relation at every time. The
animation illustrates the theorem; the derivatives in the proof establish it.

\section{Change the phase, keep the period}
Phase changes the starting point without changing the angular speed.
Indeed,
\[
  u_\phi(t)=u_0\!\left(t+\frac{\phi}{\omega}\right).
\]
Thus increasing the phase shifts the same waveform horizontally in time.
For example, $u_0(0)=0$, whereas $u_{\pi/2}(0)=1$. Both functions still
return to their starting values after time $T$.

\LemniroScene{circle-and-sine}

The three copies above have independent playback and time controls. Set two
copies to times separated by $T/4$ and compare their heights: that time offset
corresponds to a phase difference of $\pi/2$. Changing one copy's time or
playback speed does not change another copy's playback state.
The PDF retains static illustrations and this explanation; the web edition
adds the controls.

\begin{exercise}
For $v(t)=A\sin(\omega t+\phi)$, where $A\in\R$, verify that
$v''+\omega^2v=0$. Then find the values of $v(0)$ and $v'(0)$, and explain
how amplitude and phase affect these initial data.
\end{exercise}
\end{document}
